% ===========================================================================
%  E13 — Equipes na missão
% ===========================================================================
\chapter{E13 \textperiodcentered{} Equipes na missão}
\label{ep:e13}

\begin{fichaepico}
\begin{campos}
  \item[Objetivo] Quando os estudantes dividem um computador, o professor
    organiza a turma em equipes, e a plataforma passa a saber quem fez a
    missão junto.
  \item[Para quem] Professor (monta as equipes e lê a turma) e Estudante (vê a
    própria equipe). A Coordenação consulta, sem mexer.
  \item[Valor] A maioria das escolas não tem um computador para cada
    estudante. Na sala com computador dividido, quem entra na missão
    representa a equipe --- e, sem saber disso, a plataforma conta a sala como
    se cada criança tivesse a própria tela: \enquote{12 não começaram}, com os
    12 sentados ao lado de quem começou.
  \item[Estado] \parcial{} --- montar, mostrar e contar por computador estão
    prontos e em revisão; falta o Instanooria formar equipes do mesmo lado do
    feed.
\end{campos}
\end{fichaepico}

\begin{tabularx}{\linewidth}{@{}L{16mm}X L{20mm}@{}}
\toprule
\headrow \thd{ID} & \thd{Feature} & \thd{Estado}\\
\midrule
\id{F13.1} & Montar as equipes & \pronto\\
\zebra \id{F13.2} & O aluno sabe a sua equipe & \pronto\\
\id{F13.3} & Acompanhar por computador & \pronto\\
\zebra \id{F13.4} & Instanooria: equipe de lado igual & \afazer\\
\bottomrule
\end{tabularx}

\section{O que é uma equipe aqui}

\textbf{Equipe é quem divide um computador numa missão}: dois ou três
estudantes diante da mesma tela, com o login de um deles. Não é grupo
pedagógico nem equipe fixa da turma --- vale para uma missão numa turma, e a
aula seguinte pode ter outras equipes. Também não se confunde com as equipes de
resgate da Missão Tempestade, que são parte do jogo.

\begin{nota}
\textbf{Pronto, em revisão.} As histórias marcadas \pronto{} neste épico
funcionam de ponta a ponta e têm o teste que percorre a tela, mas em 25 de
setembro de 2026 ainda estavam em revisão pelo time, antes de entrar na
plataforma, e aguardam a homologação. Se a revisão pedir ajuste, a história
volta a \parcial{}, com a pendência escrita.
\end{nota}

\section{O que fica fora desta onda}

\begin{decisao}
\textbf{Decidido em 25 de setembro de 2026.}
\begin{itemize}
  \item \textbf{Só três missões aceitam equipe}: Instanooria, Missão
    Tempestade e UnooTube. Nas outras, a ficha nem mostra a seção de equipes.
  \item \textbf{A Vila do Lume fica fora}: o retrato de valores que ela
    produz é uma medida de cada criança, e é ele que define o lado de largada
    no Instanooria (\id{H06.6.3}). Duas crianças numa tela dariam um retrato
    que não é de nenhuma das duas.
  \item \textbf{O Eco no Grupo fica fora}: A Seleção o substitui
    (\id{H06.10.2}) e já nasce com equipes.
  \item \textbf{O Que a Foto Conta, Quem mandou isso?, O Rastro do Boato e O
    Manifesto da IA ficam fora}, porque serão retirados da trilha.
  \item \textbf{O atrasado não entra numa equipe automaticamente}: numa sala
    com computador dividido, quem chega atrasado senta num computador que já
    entrou na missão. Se o professor quiser registrar, põe o estudante na
    equipe à mão (\id{H13.1.3}).
  \item \textbf{A evidência pedagógica por habilidade da BNCC, valendo para
    todos da equipe}, é outro épico: \id{E14} Evidência pedagógica
    (p.~\pageref{ep:e14}), com a base pronta e à espera da primeira missão com
    régua.
\end{itemize}
\end{decisao}

% ---------------------------------------------------------------------------
\feature{F13.1}{Montar as equipes}{\pronto}

As equipes se montam na ficha da missão, aberta pela sala (\enquote{Conduzir
com esta turma}), na seção \enquote{Equipes desta turma}: \enquote{Para quando
alunos dividem um computador: quem entra na missão representa a equipe.}

\begin{historia}{H13.1.1}{Dizer quantos computadores há e quem faltou}{\pronto}
\hu{Professor}{dizer quantos computadores a turma vai usar e quem faltou
hoje}{que as equipes sejam feitas só com quem está na sala}
\aceite
\begin{itemize}
  \item \enquote{Quantos computadores vocês vão usar?} aceita de 1 até o
    número de estudantes, e logo abaixo a tela diz o que isso dá (ex.:
    \enquote{25 estudantes em 10 computadores: equipes de 2 ou 3.}).
  \item Com um computador para cada um, a tela diz \enquote{cada um no seu}
    --- e isso também vale.
  \item \enquote{Quem faltou hoje?} abre a lista da turma: quem é marcado fica
    fora do sorteio e aparece em \enquote{Sem equipe}.
  \item Marcar a falta aqui \textbf{não é chamada}: não vale como registro de
    frequência e serve só para montar as equipes desta missão.
\end{itemize}
\end{historia}

\begin{historia}{H13.1.2}{Escolher entre sortear e separar eu mesmo}{\pronto}
\hu{Professor}{escolher, a cada vez, se a plataforma sorteia as equipes ou se
eu mesmo separo}{usar o que eu sei da minha turma quando for preciso}
\aceite
\begin{itemize}
  \item A plataforma \textbf{sempre pergunta} \enquote{Como montar as
    equipes?}, com \enquote{Sortear as equipes} e \enquote{Eu separo}, sem
    nenhuma opção já marcada.
  \item \enquote{Sortear as equipes} divide quem está presente em equipes de
    tamanho parecido, e dá para mudar quem fica onde antes de salvar;
    \enquote{Sortear de novo} refaz o sorteio.
  \item Os tamanhos diferem em no máximo um: 7 estudantes em 3 computadores
    viram equipes de 3, 2 e 2 --- nunca 3, 3 e 1, com alguém sobrando
    sozinho.
  \item \enquote{Eu separo}: \enquote{A plataforma cria 10 equipes vazias e
    você põe cada estudante numa delas} --- uma equipe vazia por computador.
\end{itemize}
\end{historia}

\begin{historia}{H13.1.3}{Ajustar as equipes à mão}{\pronto}
\hu{Professor}{mudar as equipes antes de salvar}{separar quem não trabalha bem
junto, ou juntar quem precisa de ajuda}
\aceite
\begin{itemize}
  \item Posso mudar cada estudante de equipe, renomear uma equipe, criar outra
    (\enquote{Nova equipe}) e remover uma equipe.
  \item Quem ainda não está em nenhuma equipe aparece em \enquote{Sem equipe};
    com todos distribuídos, a tela diz \enquote{Todo mundo está numa equipe.}
  \item Quem foi marcado como falta continua escolhível: o atrasado que chegar
    pode ser posto numa equipe à mão.
  \item A plataforma não deixa salvar equipe vazia (\enquote{Há uma equipe sem
    ninguém. Ponha alguém nela ou remova a equipe.}), equipe sem nome
    (\enquote{Toda equipe precisa de um nome.}) nem duas equipes com o mesmo
    nome.
\end{itemize}
\end{historia}

\begin{historia}{H13.1.4}{Salvar, ver na sala e desfazer}{\pronto}
\hu{Professor}{salvar as equipes e saber, pela sala, que elas estão
prontas}{chegar na aula com a turma já organizada}
\aceite
\begin{itemize}
  \item Ao salvar, a tela confirma \enquote{Equipes salvas.} e mostra um
    cartão por equipe.
  \item Na sala da turma, a missão ganha a etiqueta \enquote{5 equipes
    prontas} --- só informa, não é um atalho.
  \item Cada estudante fica em uma equipe só por missão.
  \item Se alguém salvou equipes para a turma enquanto eu editava, a tela
    avisa \enquote{Estas equipes mudaram em outro lugar} e oferece
    \enquote{Recarregar}: nada é sobrescrito sem que eu veja.
  \item \enquote{Desfazer equipes} pede confirmação (\enquote{Desfazer as
    equipes?}): as equipes desta missão nesta turma deixam de existir, e dá
    para montar de novo quando quiser.
  \item As equipes continuam editáveis mesmo depois de a missão ser
    encerrada.
\end{itemize}
\end{historia}

\begin{historia}{H13.1.5}{Repetir as equipes da última missão}{\pronto}
\hu{Professor}{reaproveitar as equipes que a turma já usou}{não montar tudo de
novo a cada missão}
\aceite
\begin{itemize}
  \item Quando a turma já tem equipes salvas em outra missão, aparece uma
    terceira opção: \enquote{Manter as equipes da última missão}, que diz de
    qual missão elas vêm e quando foram salvas.
  \item Quem faltou hoje sai da equipe e fica em \enquote{Sem equipe}; quem saiu
    da turma desde então não volta, e a tela diz quantos são.
  \item Nada fica valendo antes de eu salvar: posso ajustar como em
    \id{H13.1.3}.
\end{itemize}
\end{historia}

\begin{historia}{H13.1.6}{Mostrar as equipes para a turma}{\pronto}
\hu{Professor}{projetar as equipes para a sala}{que cada estudante ache o
computador dele sem eu ler a lista em voz alta}
\aceite
\begin{itemize}
  \item \enquote{Mostrar equipes} abre uma janela em tela cheia com o nome de
    cada equipe e o \textbf{primeiro nome} de cada estudante.
  \item Nada de faltas, de desempenho nem de quem ficou \enquote{Sem equipe}:
    é uma tela que a turma inteira vê.
  \item Quando dois estudantes da turma têm o mesmo primeiro nome, entra a
    inicial do sobrenome (\enquote{Maria S.}). Se a inicial também empata,
    entra o último sobrenome inteiro (\enquote{Gabriela Cardoso} e
    \enquote{Gabriela Correia}): cada estudante acha o próprio nome, e o nome
    completo só aparece quando nada mais curto separa os dois.
\end{itemize}
\end{historia}

\begin{historia}{H13.1.7}{Consultar as equipes, sem mexer}{\pronto}
\hu{Coordenação da escola}{ver as equipes que o professor montou}{acompanhar a
aula sem interferir nela}
\aceite
\begin{itemize}
  \item A Coordenação vê as equipes salvas, ou \enquote{Nenhuma equipe montada
    para esta turma.}
  \item Não há botão de montar, editar nem desfazer: só quem dá aula na turma
    monta as equipes.
\end{itemize}
\end{historia}

% ---------------------------------------------------------------------------
\feature{F13.2}{O aluno sabe a sua equipe}{\pronto}

\begin{historia}{H13.2.1}{Ver com quem divido o computador}{\pronto}
\hu{Estudante}{ver, ao abrir a missão, qual é a minha equipe e com quem eu
sento}{achar o meu lugar sem depender do quadro}
\aceite
\begin{itemize}
  \item Uma faixa de uma linha, acima da missão, diz \enquote{Sua equipe:
    Equipe 3 \textperiodcentered{} com Ana e Bruno}.
  \item Os colegas aparecem só pelo primeiro nome.
  \item Quando o primeiro nome repete na equipe (o meu inclusive), entra a
    inicial do sobrenome (\enquote{Gabriela C.}); se a inicial também empata,
    o último sobrenome inteiro (\enquote{Gabriela Cardoso}). É a mesma regra
    da tela projetada (\id{H13.1.6}): o nome completo só aparece quando nada
    mais curto separa os dois.
  \item A faixa é da plataforma, fora do aplicativo imitado: a missão em si não
    muda.
  \item Sem equipe nesta missão, a faixa não aparece e a missão abre como
    sempre.
\end{itemize}
\end{historia}

% ---------------------------------------------------------------------------
\feature{F13.3}{Acompanhar por computador}{\pronto}

\begin{historia}{H13.3.1}{Ler a turma por computador}{\pronto}
\hu{Professor}{que os resultados do UnooTube e da Missão Tempestade contem
computadores quando a turma está em equipes}{saber quem de fato não começou,
em vez de ver a sala inteira como atrasada}
\aceite
\begin{itemize}
  \item Com equipes, cada computador é uma linha: a coluna passa a se chamar
    \enquote{Estudante ou equipe} e mostra o nome da equipe, quem senta nela e
    quem entrou.
  \item Os totais contam computadores (\enquote{Computadores na turma}, no
    lugar de \enquote{Estudantes na turma}), com o aviso \enquote{Contado por
    computador: cada equipe desta turma é uma linha, com os dados de quem
    entrou na missão.}
  \item Se mais de um da equipe entrou, a linha mostra quem entrou por último.
  \item Sem equipes, os resultados continuam como antes: um estudante por
    linha.
\end{itemize}
\end{historia}

\begin{historia}{H13.3.2}{Guias que explicam a aula em equipe}{\pronto}
\hu{Professor}{que o guia da aula me diga como conduzir a missão com
computador dividido}{não descobrir as regras no meio da aula}
\aceite
\begin{itemize}
  \item Os guias do UnooTube e da Missão Tempestade apontam para
    \enquote{Equipes desta turma}, na ficha da missão (Minhas turmas \seta{} a
    turma \seta{} Conduzir com esta turma).
  \item O da Missão Tempestade avisa que as 3 perguntas e o limite diário de IA
    são da conta de quem entrou: a equipe negocia em voz alta o que perguntar,
    e o que for escrito fica no nome de quem entrou.
  \item O do UnooTube lembra que os 3 vídeos salvos são da equipe.
\end{itemize}
\end{historia}

% ---------------------------------------------------------------------------
\feature{F13.4}{Instanooria: equipe de lado igual}{\afazer}

No Instanooria, cada estudante começa num lado do feed (\id{H06.6.3}). Hoje o
professor já consegue montar equipes para essa missão, mas o sorteio ainda não
olha o lado: uma equipe pode juntar estudantes dos dois lados, e quem entra
recebe o próprio lado. Esta é a próxima frente do épico.

\begin{historia}{H13.4.1}{Sortear equipes do mesmo lado do feed}{\afazer}
\hu{Professor}{que o sorteio junte na mesma equipe estudantes do mesmo lado do
feed}{que cada computador viva uma bolha só}
\aceiteproposta
\begin{itemize}
  \item Quem tem perfil da Vila do Lume vai para uma equipe do seu lado.
  \item Quem não tem perfil é \enquote{neutro} e é sorteado para completar as
    equipes.
  \item Quando a conta não fecha, a tela diz quantas equipes ficaram mistas.
  \item Os tamanhos continuam diferindo em no máximo um (\id{H13.1.2}).
\end{itemize}
\end{historia}

\begin{historia}{H13.4.2}{Quem entra recebe o lado da equipe}{\afazer}
\hu{Professor}{que qualquer um da equipe que entre na missão receba o mesmo
lado do feed}{que a bolha da equipe não dependa de quem pegou o teclado}
\aceiteproposta
\begin{itemize}
  \item O estudante que entra recebe o lado da equipe, e não o próprio.
  \item Se outro da equipe entrar depois, recebe o mesmo lado.
\end{itemize}
\end{historia}

\begin{historia}{H13.4.3}{Sem redistribuir o lado quando há equipe}{\afazer}
\hu{Professor}{que a turma com equipes não tenha o lado redistribuído}{que a
redistribuição não separe o lado de quem divide a mesma tela}
\aceiteproposta
\begin{itemize}
  \item Com equipes montadas, \enquote{Redistribuir a turma} fica desligado, e
    a tela diz por quê: o lado passa a ser da equipe.
  \item Sem equipes, a redistribuição funciona como hoje (\id{H06.6.3}).
\end{itemize}
\end{historia}

\begin{historia}{H13.4.4}{Ver as equipes no monitor da aula}{\afazer}
\hu{Professor}{ver as equipes no monitor do Instanooria}{acompanhar cada
computador ao vivo, como já faço no UnooTube e na Missão Tempestade}
\aceiteproposta
\begin{itemize}
  \item O monitor mostra o nome de cada equipe e quem senta nela.
  \item Sem equipes, o monitor continua como hoje (\id{H06.6.6}).
\end{itemize}
\end{historia}
